\section{Perpendicular standing spin waves in films with interfacial Dzyaloshinskii-Moriya interaction - \href{https://doi.org/10.1103/PhysRevB.111.064432}{PRB, 111, 064432 (2025)}}

\section*{Main points from the Paper}
\subsection*{\underline{I. Introduction:}}
\begin{enumerate}
    \item As Dzyaloshinski Moriya interaction (DMI) is known to induce the nonreciprocity in spin wave propagation within the plane of magnetic thin films.
    \item Analytical expressions are developed to estimated the effecets of DMI on spin waves by assuming the magnetization dynamics are uniform throughout a film thickness.
    \item The calculations done here show that the PSSW mode frequencies depend nonlinearly on the interfacial DMI strength and on the inverse thicknesses of the magnetic films.
    \item DMI is an antisymmetric exchange interaction caused by spin-orbit coupling and broken symmetry that enforces a preferred rotational sense in magnetic textures and their stabilization. It exists when the following conditions are met:
    \begin{enumerate}[\ding{212}]
        \item Strong spin-orbit coupling (usually from heavy metal).
        \item Broken inversion symmetry (typically at an interface).
    \end{enumerate}
    \item Generally Spin waves are the fundamental excitations of magnetic materials. 
    \item Spin waves in thin films are affected by interfacial DMI, leading to exciting phenomena such as nonreciprocal propagation, spin wave focusing, and even the loss of standing wave modes[\textcolor{red}{ref. 17-19}].
    \item Spin waves do provide the tool to measure the DMI experimentally \cite{stashkevich2015experimental,nembach2015linear}.
    \item Usually, in the theoretical studies of spin waves affected by DMI have assumed that the magnetization is uniform through the magnetic film's thickness \cite{moon2013spin,cortes2013influence}.
    \item The atomic layer method has been used to treat antiferromagnetically coupled stacks of thin films, exchange spring bilayers, and thin films with magnetoelectric coupling at interfaces \cite{moore2014spin}.
\end{enumerate}
\subsection*{\underline{II. Results and Discussion:}}
\subsubsection*{\centering A. One FM material}
\begin{enumerate}
    \item In the case of single FM material, it is assumed that the magnetization precession is uniform through the film.
    \item The full dipolar contribution is also included in the dispersion relation in Ref. \cite{moon2013spin}.
    \item The interfacial atomic DMI constant $D_{at}=4.2$ mJ$/m^2$ acts at the top layer only, which corresponds to an effective DMI used in the analytic expression of $D_{\mathrm{eff}}=D_{\mathrm{at}}/N=1.05$ mJ$/m^2$.
    \item The presence of DMI results in a pinning of a PSSW modes at the top surface, which is undetectable. The pinning is different for spin waves depending on the sign of $k_x$.
    \item One can see that without DMI, the mode is truly uniform through the film thickness, but with DMI and depending on the sign of $k_x$, the pinning results in nonuniform precession.
    \item Higher order standing spin waves are extremely sensitive to DMI because their motion is strongly confined across the film thickness, causing DMI and exchange to interact in a complex, non-additive way.
    \item It is well studied that increasing the thickness $L$ of a film will decrease the exchange energy and will therefore decrease the frequencies of higher-ordered PSSW [\textcolor{red}{ref. 36 \& 37}].
    \item The interface modifies the mode which then modifies the exchange interactions. And as a result, the frequency and damping is also modified. In our case this explains the non-monotonic damping trends.
    \item \hl{The system transitions from uniform-mode dynamics to exchange-dominated multi mode dynamics with increasing thickness}.
\end{enumerate}
\subsubsection*{\centering B. Two FM material}
\begin{enumerate}
    \item Interface position matters more than the interface strength. MgO interface is on one side only which introduces the asymmetry. Boundary conditions produce mode reshaping similar to FM bilayers.
    \item MgO in the case of my system can modify the local exchange interactions, orbital hybridization, and magnetoelastic coupling. Hence, the system behaves like a soft-hard exchange bilayer, even if chemically single phase. This justifies clearly the thickness dependent PSSW appearance in my work. [\textcolor{red}{V. Imp Line}]
    \item Thin films \ding{213} interface dominates \ding{213} uniform mode only.
    \item Thick films \ding{213} exchange resolved \ding{213} asymmetric standing modes appear.
    \item Novelty can be claimed on following points;
    \begin{enumerate}[\ding{212}]
        \item Mode-dependent damping
        \item Thickness-dependent transition from uniform to exchange-dominated dynamics.
        \item Interface-controlled PSSW emergence without explicit DMI.
    \end{enumerate}
\end{enumerate}
\begin{physicsmodule}{Related to my Paper}{Brown}
  \begin{lawbox}
    \subsection*{\sffamily\bfseries CoFe/MgO}
    \ding{212} In thicker CoFe films, perpendicular standing spin-wave modes emerge due to enhanced exchange confinement across the film thickness. These higher-order modes exhibit strong sensitivity to interfacial effects arising from the CoFe/MgO interface. Although such interfacial contributions are localized, they modify the standing-wave profile, leading to a nonlinear coupling between exchange-dominated dynamics and interface-induced fields. As a result, higher-order PSSW frequencies and damping respond more strongly than the uniform mode, consistent with the emergence of mode-dependent dynamic behavior at larger thicknesses.

    \ding{212} For CoFe/MgO, MgO modifies the exchange near the interface, and local anisotropy. This directly affects PSSW formation and damping.

    \ding{212} Studies of multilayer ferromagnetic systems have shown that perpendicular standing spin-wave modes become strongly asymmetric when magnetic properties vary across the film thickness. In such cases, interfacial effects modify the standing-wave profile rather than contributing as simple additive fields, leading to a nonlinear coupling between exchange-dominated dynamics and interface-localized interactions. As a result, higher-order PSSW modes exhibit enhanced sensitivity to interfacial asymmetry, making them effective probes of exchange nonuniformity and boundary conditions. Similar considerations apply to single-layer systems with asymmetric interfaces, where the emergence of higher-order modes reflects a transition from uniform-mode to exchange-resolved dynamics with increasing thickness.
  \end{lawbox}
\end{physicsmodule}
% \subsection*{\underline{III. Conclusion:}}